\documentclass[a4paper, 12pt]{article}
\usepackage[T2A]{fontenc}
\usepackage[english,russian]{babel}
\usepackage{amsmath}
\usepackage{amsthm}
\usepackage{amsfonts}
\usepackage{amssymb}
\usepackage{pgfplots}

\usepackage{enumitem}
% Абзацы и геометрия страницы
\usepackage{indentfirst}
%\usepackage{geometry}
\usepackage[
    top=2cm,
    bottom=2cm,
    left=1cm,
    right=1cm
]{geometry}
\begin{document}


\begin{center}
    МИНИСТЕРСТВО НАУКИ И ВЫСШЕГО ОБРАЗОВАНИЯ РОССИЙСКОЙ ФЕДЕРАЦИИ

    ГОСУДАРСТВЕННОЕ ОБРАЗОВАТЕЛЬНОЕ УЧРЕЖДЕНИЕ ВЫСШЕГО ПРОФЕССИОНАЛЬНОГО ОБРАЗОВАНИЯ

    РЯЗАНСКИЙ ГОСУДАРСТВЕННЫЙ РАДИОТЕХНИЧЕСКИЙ УНИВЕРСИТЕТ

    ИМ. В. Ф. УТКИНА

    КАФЕДРА ВЫЧИСЛИТЕЛЬНОЙ И ПРИКЛАДНОЙ МАТЕМАТИКИ
\end{center}

\vspace{5 cm}

\begin{center}
    ЛАБОРАТОРНАЯ РАБОТА № 5

    ПО ДИСЦИПЛИНЕ "ТЕОРЕТИЧЕСКАЯ ИНФОРМАТИКА"

    НА ТЕМУ "ОПЕРАЦИИ С ЧИСЛАМИ С ФИКСИРОВАННОЙ ТОЧКОЙ"

    ВАРИАНТ 14

\end{center}

\vspace{7 cm}

\begin{flushright}
    Выполнил:

    Студент группы 6413

    Косяков Д.А.

    Проверил

    Доцент кафедры ВПМ

    Никичкин Борис Владимирович
\end{flushright}

\begin{center}
    РЯЗАНЬ 2026
\end{center}

\pagebreak

\textbf{ Цель работы:} выполнить сложение чисел с фиксированной точкой, представленных в прямом коде в двоичной системе счисления.

\textbf{ Задание:} перевести числа А и В в десятичную систему счисления,
сложить их и результат записать в двоичной системе счисления.
Числа А и В представлены в прямом коде.

\textbf{ Выполнение задания:} заданы $A_2$ = 100000,1111 и $B_2$ = 110011,0111 (прямой код).

В прямом коде старший (самый левый) разряд --- знаковый: 0 означает, что число положительное,
1 --- что число отрицательное. Остальные разряды хранят модуль числа.

Разделим число А на знак и модуль: знак --- 1, модуль --- 00000,1111. Знак равен 1, значит число отрицательное.
Переведем модуль числа А в десятичную систему счисления. Разряды целой части умножаем на
положительные степени двойки, разряды дробной части --- на отрицательные:
$$|A_{10}| = 0 \cdot 2^4 + 0 \cdot 2^3 + 0 \cdot 2^2 + 0 \cdot 2^1 + 0 \cdot 2^0
+ 1 \cdot 2^{-1} + 1 \cdot 2^{-2} + 1 \cdot 2^{-3} + 1 \cdot 2^{-4}$$
$$|A_{10}| = 0{,}5 + 0{,}25 + 0{,}125 + 0{,}0625 = 0{,}9375$$
$$A_{10} = -0{,}9375$$

Разделим число В на знак и модуль: знак --- 1, модуль --- 10011,0111. Знак равен 1, значит число отрицательное.
Переведем модуль числа В в десятичную систему счисления:
$$|B_{10}| = 1 \cdot 2^4 + 0 \cdot 2^3 + 0 \cdot 2^2 + 1 \cdot 2^1 + 1 \cdot 2^0
+ 0 \cdot 2^{-1} + 1 \cdot 2^{-2} + 1 \cdot 2^{-3} + 1 \cdot 2^{-4}$$
$$|B_{10}| = 16 + 2 + 1 + 0{,}25 + 0{,}125 + 0{,}0625 = 19{,}4375$$
$$B_{10} = -19{,}4375$$

Вычисляем сумму чисел $A_{10}$ и $B_{10}$:
$$A_{10} + B_{10} = -0{,}9375 + (-19{,}4375) = -20{,}375$$

Переведем модуль суммы в двоичную систему счисления. Целую часть переводим делением на 2:

$$20 : 2 = 10, \text{ остаток } 0$$
$$10 : 2 = 5, \text{ остаток } 0$$
$$5 : 2 = 2, \text{ остаток } 1$$
$$2 : 2 = 1, \text{ остаток } 0$$
$$1 : 2 = 0, \text{ остаток } 1$$

Получаем $20_{10} = 10100_2$.

Дробную часть переводим умножением на 2:

$$0{,}375 \cdot 2 = 0{,}75, \text{ целая часть } 0$$
$$0{,}75 \cdot 2 = 1{,}5, \text{ целая часть } 1$$
$$0{,}5 \cdot 2 = 1{,}0, \text{ целая часть } 1$$

Получаем $0{,}375_{10} = 0{,}011_2$.

Значит $20{,}375_{10} = 10100{,}011_2$, а $-20{,}375_{10} = -10100{,}011_2$.

В прямом коде (как у исходных чисел: 1 знаковый разряд, 5 разрядов целой части и 4 разряда дробной части)
ставим в знаковый разряд 1: $-20{,}375_{10} = 110100{,}0110_2$ (прямой код).

Проверка: числа одного знака в прямом коде складываются так --- модули складываются, знак сохраняется:
$00000{,}1111_2 + 10011{,}0111_2 = 10100{,}0110_2$, знак --- минус. Результат совпадает.
\vspace{1cm}

\textbf{Отчет:}
\begin{enumerate}[nosep]
    \item А = 100000,1111, В = 110011,0111 (прямой код).
    \item $A_{10} = -0{,}9375$, $B_{10} = -19{,}4375$.
    \item $A_{10} + B_{10} = -20{,}375$.
    \item $-20{,}375_{10} = -10100{,}011_2$, в прямом коде 110100,0110.
\end{enumerate}

\end{document}
